\documentclass[]{article}

%opening
\title{Saddle Thrombus Yogurt}
\author{}

\begin{document}

\maketitle

\begin{abstract}

\end{abstract}

\section{Introduction}

Saddle thrombus, also known as feline aortic thromboembolism (FATE), is a serious and often life-threatening condition in cats where a blood clot (thrombus) forms, typically originating from the heart, and travels through the bloodstream until it lodges at the point where the aorta bifurcates into the iliac arteries that supply blood to the hind legs. This blockage, which resembles a "saddle" across the aortic split, severely restricts or cuts off blood flow to the rear limbs, leading to sudden onset of symptoms such as paralysis, severe pain, cold paws, and absent pulses in the hind legs. It is most commonly associated with underlying cardiac diseases, particularly hypertrophic cardiomyopathy (HCM), which affects the heart's ability to pump blood effectively and promotes clot formation. Other risk factors include hyperthyroidism or certain infections, but heart disease remains the primary driver in the majority of cases. 

\section{Standard of Care}
The conventional veterinary standard of care for saddle thrombus focuses on rapid intervention to manage pain, stabilize the cat, and address both the clot and any underlying conditions, though the prognosis is generally guarded to poor due to high recurrence rates and complications. Upon presentation, immediate priorities include aggressive pain relief using opioids (e.g., buprenorphine or fentanyl) and sedation if needed, alongside supportive measures like intravenous fluids to maintain hydration and blood pressure, oxygen therapy, and warming the affected limbs to improve circulation. Anticoagulant medications, such as low-molecular-weight heparin or antiplatelet drugs like clopidogrel, are routinely administered to prevent further clot formation or extension. 

Treatment of the underlying heart disease is essential, often involving medications like beta-blockers, ACE inhibitors, or pimobendan to manage cardiomyopathy. In some cases, thrombolytic therapy (e.g., using drugs like alteplase to dissolve the clot) or surgical embolectomy may be attempted, but these are not standard first-line options due to significant risks, including bleeding or reperfusion injury, and are typically reserved for specialized veterinary centers. Physical therapy and nursing care, such as bladder expression and preventing pressure sores, are also part of ongoing management for cats that survive the acute phase. Unfortunately, due to the intense pain and low long-term survival rates (often less than 50\% even with treatment), humane euthanasia is frequently discussed as a compassionate option if the cat does not respond quickly or if recurrence is likely. 

Early diagnosis via physical exam, ultrasound, and blood tests is crucial, and cat owners should seek emergency veterinary care immediately upon noticing symptoms.


\end{document}
